\chapter{动机与去动机}\label{s:motivation}

学习者需要鼓励才会踏进陌生的领域，
所以本章讨论老师可以怎样激励他们。
更重要的是，
它讲老师会怎样\emph{打击}他们，
以及如何避免这么做。

我们的出发点是
\gref{g:extrinsic-motivation}{外在动机}
和\gref{g:intrinsic-motivation}{内在动机}之间的区别。
前者是我们为了免于惩罚或获得奖赏而做事时感到的，
后者是我们发现某件事能让自己感到充实时感到的。
两者都影响大多数情境——例如，
人们教书既因为喜欢、也因为能拿到钱——
但我们在受内在动机驱动时学得最好~\cite{Wlod2017}。
根据\hreffoot{https://en.wikipedia.org/wiki/Self-determination\_theory}{自我决定论}，
内在动机的三个驱动力是：

\begin{description}

\item[胜任感：]
  觉得自己知道自己在做什么。\index{动机!胜任感}

\item[自主感：]
  觉得能掌控自己的命运。\index{动机!自主感}

\item[关联感：]
  觉得和他人有连接。\index{动机!关联感}

\end{description}

一门设计得好的课会同时鼓励这三样。
例如，
一道编程练习可以
让学习者练习解决一个更大问题所需的工具（胜任感）、
让他们按自己想要的任何顺序处理该问题的各个部分（自主感）、
并允许他们和同伴交流（关联感）。

\begin{aside}{分数的麻烦}
  我这辈子从没有过听众。我的听众是一份评分标准。\\
  —— 引自\hreffoot{https://twitter.com/figuralities/status/987330064571387906}{Matt Tierney}

  分数以及它们对学习的扭曲，常被当作外在动机的例子，
  但正如~\cite{Mill2016a}所观察到的，
  它们短期内不会消失，
  所以想建一套忽略它们的制度是没意义的。
  相反，~\cite{Lang2013}探讨了强调分数的课程
  如何会诱使学习者作弊，
  并就如何减弱这种效应给出了一些建议；
  而~\cite{Covi2017}考察了在院校教育中
  平衡内在和外在动机这个更大的问题，
  \cite{Bigg2011}所提倡的\hreffoot{https://en.wikipedia.org/wiki/Constructive\_alignment}{一致性建构}方法
  则设法让学习活动和学习成效彼此对齐。
\end{aside}

\cite{Ambr2010}包含一份基于证据的激励学习者方法清单。
它们都不出人意料——很难想象
有人会说我们\emph{不该}识别并奖励我们所看重的东西——
但用它来检查课程、确保课程至少做到了其中几样，是有用的。
我特别喜欢的一个策略是，
让那些挣扎过但最终成功的学习者
进来把他们的故事讲给全班听。
学习者更可能相信和自己相似的人讲的故事~\cite{Mill2016a}，
而走过你课程的人，
总能给出你永远想不到的建议。

\begin{aside}{不只是为了学习者}
  教育中关于动机的讨论，常常忽略也需要激励\emph{老师}。\index{动机!老师的}
  学习者会对老师的热情做出回应，
  而老师（尤其是志愿者）需要在乎一个主题，才能一直教下去。
  这是协同教学（\secref{s:classroom-together}）的又一个有力理由：
  就像有个跑步搭子会让你更可能坚持跑下去，
  有个教学搭档，能帮你在那些
  你感冒了、
  投影灯泡烧了、
  谁都不知道替换灯泡在哪儿、
  而且说真的，
  他们\emph{又}在施工了吗？
  的日子里，照样爬起来、出门去上课。
\end{aside}

老师还可以做其他一些正面的事。
\cite{Bark2014}发现了驱动所有学习者留存的三个因素：
有意义的作业、
师生互动、
以及学习者就作业展开协作。
相对于期望的进度和工作量也是显著的驱动因素，
但主要对男性学习者而言。
\emph{不}驱动留存的
是与助教的互动，
以及课外活动中与同伴的互动。
这些结果似乎显而易见，
但反过来也会显得显而易见：
如果研究发现课外活动确实驱动了留存，
我们也会觉得这说得通。
值得注意的是，
四个留存驱动因素中的两个（师生互动和学习者协作）
在网上复现要花更多功夫（\chapref{s:online}）。

\seclbl{真实任务}{s:motivation-authentic}

正如 Dylan Wiliam 在~\cite{Hend2017}中指出的，\index{Wiliam, Dylan}
动机并不总能导向成就，
但成就几乎总能导向动机：
学习者从自己的成功中得到的动机，远多于被告知他们有多棒。
我们可以在教学中利用这个想法，
办法是画一张网格，横纵轴分别是「掌握所需的平均时间」
和「掌握后的有用程度」（\figref{f:motivation-what}）。

\figpdf{figures/what-to-teach.pdf}{教什么}{f:motivation-what}{}

学起来快、又马上有用的东西应该先教，
即便它们不被已经胜任的人认为有多根本，
因为早期的几次小成功会建立学习者对自己和老师的信心。
反过来，
难学、而且对处在当前发展阶段的学习者没用的东西，
应该干脆跳过；
而对角线上的主题则需要互相权衡。

\newpage
\begin{aside}{对谁有用？}
  如果某人想建网站，
  像递归和可计算性这样的基础计算机科学概念，
  可能就落在这张网格的右下角。
  这并不意味着它们不值得学，
  但如果我们的目标是激励人，
  它们就可以、也应该被推后。
  反过来，
  一位选修编程课来活动脑筋的老年人，
  可能更愿意探索这些大想法，而不是做任何实用的东西。
  你在做网格时，
  应该心里想着你的学习者画像
  （第~\ref{s:process-personas}~节）。
  如果不同画像把同一主题放到了很不一样的位置，
  你就应该考虑开设不同的课程。
\end{aside}

一个被充分研究的、既优先考虑有用的东西、又不牺牲根本东西的例子，
是佐治亚理工学院开发的媒体计算方法~\cite{Guzd2013}。\index{媒体计算}
学习者写的第一批程序，不是打印「hello world」或求前十个数之和，
而可能是打开一张图片、
调整大小生成一张缩略图、
并保存结果。
这是一个\gref{g:authentic-task}{真实任务}，
也就是学习者相信自己真会在现实生活里做的事。
它也有一个\gref{g:tangible-artifact}{有形造物}：
如果图片出来尺寸不对，
学习者手头就有东西可以指导自己调试。
\cite{Lee2013}描述了把这种方法从 Python 改到 MATLAB 的尝试，
而其他人则在围绕数据科学、图像处理
和生物学构建类似的课程~\cite{Dahl2018,Meys2018,Ritz2018}。

在给学习者真实问题和操练解决这些问题所需的单项技能之间，永远会有张力：
毕竟，
程序员在工作中不会去答多项选择题，
就像音乐家不会在听众面前一遍遍弹音阶一样。
找到平衡很难，
但第一步是把任何武断或无意义的东西拿掉。
例如，
编程示例不该用叫 \texttt{foo} 和 \texttt{bar} 的变量；
如果你要让学习者给一个列表排序，
就用一个歌单，而不是「aaa」和「bbb」这样的字符串。

\seclbl{去动机}{s:motivation-demotivation}

\begin{quote}

  女性离开计算领域，不是因为她们不知道那是什么样子；
  她们离开，是因为她们\emph{知道}。\\
  —— 出处不一

\end{quote}

如果你在散养式环境里教书，
你的学习者是志愿者，
而且多半想来你的课堂。
所以激励他们，不如「别打击他们」来得要紧。
遗憾的是，
你很容易无意中就打击了人。
例如，
\cite{Cher2009}报告了四项研究，表明\index{去动机!环境因素}
微妙的环境线索，会对不同性别的人对计算的兴趣产生可测量的影响：
把计算机科学教室里的物品，从那些被认为是计算机科学刻板印象的东西
（例如《星际迷航》海报和电子游戏）
换成不被认为刻板的物品（例如自然海报和电话簿），
就把女本科生的兴趣提升到了与男同学同等的水平。
同样，
\cite{Gauc2011}报告了三项研究，表明
招聘材料里常用的性别化措辞，
会维持传统上男性主导职业中的性别不平等。

成人学习者有三个主要的去动机因素：

\begin{description}

\item[不可预测]
  会打击人，因为
  如果他们所做和所达成的结果之间没有可靠的关联，
  他们就没有理由去尝试做任何事。\index{去动机!不可预测}

\item[冷漠]
  会打击人，因为
  认为老师或教育系统不在乎自己、
  也不在乎材料的学习者，也不会去在乎它。\index{去动机!冷漠}

\item[不公平]
  出于显而易见的理由，会打击处于不利地位的人。
  让人意外的是，它也会打击从不公平中受益的人：
  不管意识到与否，
  他们担心自己有朝一日也会落入不利的一方~\cite{Wilk2011}。\index{去动机!不公平}

\end{description}

在极端情况下，
学习者可能发展出\gref{g:learned-helplessness}{习得性无助}：
当在一个自己无法改变的处境里反复遭受负面反馈时，
他们可能学会连自己本可以改变的事情都不再尝试去改变。

打击学习者最快、最确凿的方式之一，
是使用某种暗示「有些人是天生的程序员、有些人不是」的语言。
Guzdial 把这叫做
\hreffoot{https://cacm.acm.org/blogs/blog-cacm/189498-top-10-myths-about-teaching-computer-science/fulltext}{关于计算机教学最大的迷思}，
而\cite{Pati2016}支持了这一点，他们表明
人们会在根本不存在的地方看到「极客基因」的证据。\index{极客基因（并不存在）}
他们分析了 778 门大学课程的成绩分布，发现只有 5.8\% 显示出
多峰分布的迹象，
也就是每二十个班里只有一个班显示出有两类明显不同的学习者群体。
然后他们向 53 位计算机科学教授展示了模棱两可的成绩分布直方图；
那些相信有些人生来更适合学计算机科学的人，
比不相信的人更可能把这些分布看成双峰的。

这些信念之所以要紧，是因为老师会照它们行动~\cite{Brop1983}。
如果一位老师相信某个学习者很可能表现出色，
他们就会自然地（往往是无意识地）把注意力放在那个学习者身上，
后者由于获得了更多关注而实现了老师的期望，
这又反过来似乎印证了老师的信念。
遗憾的是，
几乎没有迹象表明，\cite{Pati2016}所呈现的那类证据本身，
就足以打破这个恶性循环{\ldots}

下面还有一些会打击你学习者的具体做法：

\begin{description}

\item[居高临下或轻蔑的态度，]
  来自老师或某个学习者。

\item[告诉他们现有的技能是垃圾。]
  Unix 用户嘲笑 Windows，
  各种各样的程序员拿 Excel 开玩笑，
  而且不管你已经在用哪种 Web 应用框架，
  总会有程序员告诉你它过时了。
  学习者常常已经投入了大量时间和精力才获得他们现有的技能；
  贬低它们，是很有效的一招，保证
  他们不会再听你说的任何别的话。

\item[和班里最进阶的学习者一头扎进复杂或细节的技术讨论。]

\item[假装你懂的比你实际懂的更多。]
  如果你对自己的知识边界坦率，
  学习者会更信任你，
  也更可能提问和求助。

\item[用那个「就」字，或假作惊讶。]
  正如\chapref{s:memory}里讨论的，
  说「我不敢相信你连 X 都不知道」或「你从没听说过 Y？」
  这类话，
  是在向学习者发出信号：
  老师觉得他们的问题微不足道，
  所以他们一定很笨，连这都想不出来。

\item[软件安装的头痛。]
  人们第一次接触编程或新编程工具时，常常很受打击，
  而「相信某样东西很难学」会变成一个自我实现的预言。
  这不仅仅是因为配置要花时间，
  或是因为觉得「不得不调试一个恰好依赖自己还不具备的知识的东西」很不公平。
  真正的问题是，每一次这样的失败都会强化一种信念：
  如果他们照老样子做事，
  在赶下周四的截止日期时会更有胜算。

\end{description}

在网上打击人，比当面还容易，
但现在已经有了基于证据的应对策略。
\cite{Ford2016}发现，在 \hreffoot{https://stackoverflow.com/}{Stack Overflow} 上，有五个贡献障碍
被女性认为显著比男性更成问题：
不了解网站功能、
觉得自己没资格回答问题、
社区规模吓人、
和陌生人互动或依赖陌生人不自在、
以及觉得在网上搜东西不算「真正的工作」。
害怕负面反馈没能进入这个清单，
但如果作者对统计阈值不那么严格，它会是下一个被加进去的。
所有这些因素，在当面和在线两种场景下都可以、也应该得到处理，
用\secref{s:motivation-inclusivity}里的那类方法，
而且这么做对每个人都有好处~\cite{Sved2016}。

\begin{aside}{有效失败与特权}
  近期有些工作探究了\gref{g:productive-failure}{有效失败}，
  即刻意给学习者一些用他们现有知识解决不了的问题，
  让他们必须走出去获取新信息才能取得进展~\cite{Kapu2016}。
  有效失败表面上让人想起科技圈的「快速失败、经常失败」口号，
  但后者更多是特权的标志，而不是理解的标志。
  人们只有在确信自己还能再得到一次机会时，才承受得起庆祝失败；
  你的许多学习者，
  以及许多来自边缘化或弱势群体的人，
  并不能确信这一点，
  而假设「失败是一种选项」，是打击他们的一大妙招。
\end{aside}

\subsection*{冒名顶替综合症}

\grefdex{g:impostor-syndrome}{冒名顶替综合症}{冒名顶替综合症}
是一种认为自己取得成就是侥幸、
并伴随「总有一天会被人看穿」的恐惧的信念。
它在从事公开可见工作的高成就者中很常见，
但对代表性不足群体成员的影响不成比例地大：
正如\secref{s:pck-now}里讨论的，
\cite{Wilc2018}发现，
在入门编程课里，此前接触过计算的女性学习者在所有方面都胜过男同学，
却一直对自己的能力更缺乏信心，
部分原因是社会不断以微妙和不那么微妙的方式发出信号，
说她们其实并不属于这里。

传统课堂会助长冒名顶替综合症。
学业常常是独自或在小群体里进行的，
但结果却会被公开分享和批评。
结果就是，
我们很少看到别人是怎么挣扎着才完成作业的，
这又会滋养「其他所有人都觉得这很容易」的信念。
那些已经感受到额外压力、要证明自己的代表性不足群体成员，
可能受影响尤甚。

Ada Initiative 制定了一些\index{冒名顶替综合症!对抗}
\hreffoot{https://www.usenix.org/blog/impostor-syndrome-proof-yourself-and-your-community}{指南}，
用来对抗你自己的冒名顶替综合症，
其中包括：

\begin{description}

\item[和你信任的人谈谈这个问题。]
  当你从别人那里听说冒名顶替综合症是普遍问题时，
  就更难相信自己那种「我是个骗子」的感觉是真的。

\item[参加一次线下的冒名顶替综合症活动。]
  没有什么比身处一屋子你尊敬的人当中、
  却发现其中 90\% 都有冒名顶替综合症更有说服力了。

\item[留意你的用词，因为会影响你的思维。]
  说「我不是这方面的专家，但是{\ldots}」这样的话，
  会削弱你实际拥有的知识。

\item[把你的领域教给别人。]
  你会对自己的知识和技能更有信心，
  也能帮别人避开一些冒名顶替综合症的暗礁。

\item[提问。]
  如果你认为自己本该知道答案，提问可能会让你发怵，
  但得到答案能消除那种不确定、害怕失败的漫长煎熬。

\item[建立同盟。]
  去安慰并壮大你的朋友，
  他们也会反过来安慰并壮大你。
  （如果他们不这么做，你也许该考虑换些新朋友{\ldots}）

\item[认领你的成就。]
  持续主动地记录并回顾你做过什么、
  建起了什么、
  取得了哪些成功。

\end{description}

作为老师，
你可以通过分享自己犯过的错误、或自己挣扎着学会的东西，
来帮人们应对他们的冒名顶替综合症。
这会让全班安心，觉得「觉得某个主题难」是没问题的。
对群体坦诚也能建立信任，
给他们提问的信心。
（现场编程在这方面很棒：
正如\secref{s:performance-live}里指出的，
你的笔误向全班表明你也是人。）
频繁的形成性评价也有帮助，
尤其是当学习者看到你根据他们的结果调整教学内容或进度时。

\subsection*{思维模式与刻板印象威胁}

Carol Dweck 等人\index{Dweck, Carol}
研究了\gref{g:fixed-mindset}{固定型思维}
与\gref{g:growth-mindset}{成长型思维}对学习成效的差异。
如果人们相信某个领域的能力是天生的
（也就是你要么「有那个基因」，要么没有），
\emph{每个人}都会表现更差，
包括那些被认为占优势的人。
原因是，如果某人一开始没做好，
他们就会假定自己缺乏那种天赋，
这又会影响他们今后的表现。
另一方面，
如果人们相信一项技能是可以学会、可以提升的，
他们平均会做得更好。

\hreffoot{https://educhatter.wordpress.com/2017/03/26/growth-mindset-is-the-theory-flawed-or-has-gm-been-debased-in-the-classroom/}{有人担心}
成长型思维被过度吹捧，
或者把关于它的研究落到实践，要比它更热情的倡导者所暗示的难得多~\cite{Sisk2018}。
不过，
社会经济地位较低、或学业上有风险的学习者，似乎确实能从思维模式干预中受益。

另一个被广泛讨论的效应是\gref{g:stereotype-threat}{刻板印象威胁}~\cite{Stee2011}。
提醒人们注意到负面的刻板印象，
即便以微妙的方式，
也会让他们焦虑于「万一证实了那些刻板印象」的风险，
这又会降低他们的表现。
同样，
也有人担心
\hreffoot{https://www.psychologytoday.com/blog/rabble-rouser/201512/is-stereotype-threat-overcooked-overstated-and-oversold}{关键研究的可重复性}，
而且这个术语被以许多不同的方式使用，让问题更加模糊~\cite{Shap2007}，
但没有人会主张在课上提起刻板印象能帮到学习者。

\seclbl{无障碍}{s:motivation-accessibility}
\index{无障碍}

让课程和练习超出某些人够得着的范围，几乎是最能打击人的事，
而且很容易无意中做到。
例如，
我写的第一批在线编程课，在幻灯片旁边有旁白的文字稿，
却没有包含真正的源代码：
源代码在 PowerPoint 幻灯片的截图里。
所以用一个\hreffoot{https://en.wikipedia.org/wiki/Screen\_reader}{屏幕阅读器}的人
能听到别人在说这个程序什么，
却不知道这个程序到底是什么。
要照顾到每个学习者的需求并不总是可行，
但给图片加上描述性说明、
并让无法使用鼠标的人也能操作导航控件，
能带来很大差别。

\begin{aside}{路缘坡道}
  让材料变得无障碍，对每个人都有帮助，
  而不只是面对挑战的人。
  \hreffoot{https://en.wikipedia.org/wiki/Curb\_cut}{路缘坡道}——连接人行道和街道的小斜坡——
  最初是为了方便身体有障碍的人出行，
  结果对有婴儿车和购物车的人同样有用。
  同样，
  给图片加说明，不只是帮视障人士：
  它还让搜索引擎更容易找到并索引这些图片。
\end{aside}

让课程变得无障碍，第一步也是最重要的一步，
是让残障人士参与决策：
口号\emph{\hreffoot{https://en.wikipedia.org/wiki/Nothing\_About\_Us\_Without\_Us}{nihil de nobis, sine nobis}}
（字面意思是「没有我们的参与，就不要做关于我们的决定」）
早于无障碍权利运动，
但永远是正确的起点。
几条具体建议是：

\begin{description}

\item[弄清你需要做什么。]
  \hreffoot{https://accessibility.blog.gov.uk/2016/09/02/dos-and-donts-on-designing-for-accessibility/}{这些海报}中的每一张，
  都给出了针对自闭症谱系人士、
  屏幕阅读器用户、
  以及低视力、
  身体或运动障碍、
  听障、
  和阅读障碍人士的该做与不该做。

\item[不要一次全做完。]
  上一条里描述的那些改进可能显得相当吓人，
  所以一次只做一项。

\item[先做容易的。]
  字体大小、
  使用夹式麦克风让人更容易听清你、
  检查你的配色选择，都是很好的起点。

\item[知道自己做得怎么样。]
  像 \hreffoot{http://webaim.org/}{WebAIM} 这样的网站，能让你检查
  你的在线材料对视障用户的可及程度。

\end{description}

\cite{Coom2012,Burg2015}是关于无障碍视觉设计的好指南。
它们的建议包括：

\begin{description}

\item[用真正的标题和其他结构标记来排版文档，]
  而不是只改字体大小和样式。

\item[避免只用颜色在文字或图形中传达意义。]
  相反，用颜色加上不同的交叉阴影图案
  （这也能让材料在黑白打印时仍然看得懂）。

\item[去掉多余元素，]
  而不是只把它们隐藏起来，
  因为屏幕阅读器往往仍会把它们读出来。

\item[允许自定节奏和重复，]
  以照顾有阅读或听力问题的人。

\item[在视频里加上对屏幕所发生事情的旁白]
  （现场编程时边打边说）。

\end{description}

\subsection*{勺子}
\index{勺子（隐喻）}

2003 年，
Christine Miserandino 开始用\hreffoot{https://butyoudontlooksick.com/articles/written-by-christine/the-spoon-theory/}{勺子}
来解释带着慢性病生活是什么样子。
健康的人每天开始时有无限多的勺子，
但患红斑狼疮或其他消耗性疾病的人只有几把，
而且他们做的每件事都要花掉一把。
起床？
一把勺子。
做顿饭？
又一把，很快你就用光了。

\begin{quote}

  你生病时可没法随便套件衣服就出门{\ldots}
  如果我那天手疼，有扣子的衣服就免谈。
  如果我那天有淤青，
  我就得穿长袖；
  如果我在发烧，我就需要一件毛衣保暖，等等。
  如果我在掉头发，我就得花更多时间让自己看起来像样，
  然后你还得再算上五分钟，用来难过
  ——为做这些花了你两个小时而难过。

\end{quote}

正如\hreffoot{https://patitsas.blogspot.com/2018/03/spoons-are-form-of-capital.html}{Elizabeth Patitsas 所主张的}，
勺子很多的人能积攒更多勺子，
而供应有限的人可能连追赶都很吃力。
在设计课程和练习时，
记住你有些学习者可能有着不明显的身心障碍。
拿不准时，就去问：
他们几乎肯定比任何人都更了解什么管用、什么不管用。

\seclbl{包容性}{s:motivation-inclusivity}

\grefdex{g:inclusivity}{包容性}{包容性}是一种
把原本可能被排斥或边缘化的人纳入进来的方针。
在计算领域，
它意味着主动努力，对女性、
代表性不足的种族或族裔群体、
各种性取向的人、
老年人、
面对身体挑战的人、
曾经入狱的人、
经济上弱势的人，
以及所有不符合硅谷富裕白人/亚裔男性人口画像的人，更加友善。
\figref{f:motivation-women-in-cs}（来自\hreffoot{https://www.npr.org/sections/money/2014/10/21/357629765/when-women-stopped-coding}{NPR}）
用图展示了计算的排斥性文化对女性的影响。

\figimg{figures/women-coding.png}{美国计算机科学专业的女性比例}{f:motivation-women-in-cs}{}

\cite{Lee2017}是一份简短、实用、并引用了研究文献的指南。
它描述的做法能帮到属于一个或多个被边缘化或被排斥群体的学习者，
但也同样能激励其他所有人。
它们是以整学期课程的口吻写的，
但很多都可以用在一天的工作坊和其他散养式环境里：

\begin{description}

\item[请学习者在工作坊之前给你发邮件，]
  说明他们认为这项培训能怎样帮他们实现自己的目标。

\item[检查你的笔记，]
  确保它们没有性别化的代词、包含文化上多样的名字，等等。

\item[强调重要的是他们学习的速度，]
  而不是他们起步时所拥有的优势或劣势。

\item[鼓励结对编程，]
  但先演示一次，好让学习者理解驾驶员和领航员的角色。

\item[主动缓减某些学习者可能觉得吓人的行为，]
  例如使用行话，或问一些其实是为了炫耀知识的「问题」。

\end{description}

支持来自边缘化群体学习者的一种办法，
是让人们以小组、而不是个人的形式报名参加工作坊。
这样，
屋里每个人事先就知道自己会和信任的人在一起，
这增加了他们真的来参加的可能性。
它在工作坊之后也有帮助：
如果人们是和自己的朋友或同事一起来的，
他们就能一起合作，运用学到的东西。

更根本地说，
课程作者需要把每个人的整体处境都考虑进去。
例如，
\cite{DiSa2014a}发现，参加一个游戏测试项目的非裔美国男性中有 65\% 后来去学了计算，
部分原因是这个项目的游戏部分是他们同伴所看重的。
\cite{Lach2018}探究了创造包容性内容的两种一般策略、
以及与之相关的风险：

\begin{description}

\item[{\grefdex{g:community-representation}{社群表征}{社群表征}}]
  利用课后导师、或来自学习者社区的角色榜样，
  或者把社区叙事和历史作为计算项目基础的活动，
  来凸显学习者的社会身份、历史和社群网络。
  这种做法的最大风险是流于表面，
  例如用电脑做幻灯片，而不是做任何真正的计算。

\item[{\grefdex{g:computational-integration}{计算式整合}{计算式整合}}]
  把学习者社区里的想法纳入进来，
  比如在可视化编程环境里复现原住民的图形设计。
  这里的最大风险是文化挪用，
  例如使用某些做法而不说明其来源。

\end{description}

拿不准时，
就去问你的学习者和社区成员，他们认为你该怎么做。
我们在\chapref{s:community}里会回到这一点。

\begin{aside}{行为准则也是一种无障碍}
  我们在\secref{s:classroom-coc}里说过，课堂应该执行一份像\appref{s:conduct}那样的行为准则。
  这是一种无障碍：
  闭路字幕让视频对听障人士可及，
  而行为准则让课程对原本会被边缘化的人可及。
\end{aside}

\subsection*{走出缺陷模型}

取决于你相信谁的数据，
拿到计算机科学学位的人里只有 12—18\% 是女性，
这还不到 1980 年代中期所见比例的一半
（\figref{f:motivation-gender}，来自~\cite{Robe2017}）。
而西方国家在计算领域女性比例如此之低，才是异类：
在其他地方，女性往往仍占计算机科学学生的 30—40\%~\cite{Galp2002,Varm2015}。

\figimg{figures/enrollment.png}{授予的学位与女性入学人数}{f:motivation-gender}{}

既然女性不可能在过去三十年里发生了剧烈变化，
我们就得去找结构性原因，来理解哪里出了问题、以及怎么修。
一种解释是家庭电脑从 1980 年代起被营销成「男孩的玩具」~\cite{Marg2003}；
另一种是计算机科学系在 1980 年代和 2000 年代两度面对入学人数的爆炸式增长时，
通过改变录取要求来应对~\cite{Robe2017}。
这些因素对不受其影响的人来说，可能都不显得戏剧性，
但它们就像水滴石穿：
久而久之，它们侵蚀动机，也侵蚀参与度。

修好这一切的第一步也是最重要的一步，
是停止用「漏水的管道」这样的说法来思考~\cite{Mill2015}。
更一般地说，
我们需要走出\gref{g:deficit-model}{缺陷模型}，
也就是不再认为代表性不足群体的成员缺少什么、
并因此该为他们没能出头负责。
相信这一点，会把负担压到那些本就得付出额外努力才能克服结构性不公的人身上，
并且（并非巧合地）给从现有格局中受益的人一个借口，
让他们不必太仔细地审视自己。

\begin{aside}{重写历史}
  \cite{Abba2012}描述了那些塑造了计算早期历史的女性
  的职业生涯和成就，
  而她们太常被从历史中抹去；
  \cite{Ensm2003,Ensm2012}描述了编程是如何在 1960 年代从一种女性职业变成男性职业的，
  而~\cite{Hick2018}审视了英国如何因为系统地歧视其最有资格的劳动者——
  女性——而失去了在计算领域的早期主导地位。
  （三本书的评论见\cite{Milt2018}。）
  讨论这段历史会让计算领域的某些男性非常不舒服；
  在我看来，
  这正是该讲它的一个好理由。
\end{aside}

电子游戏里的厌女、
招聘中用「文化契合」来为有意或无意的偏见开脱、
围绕性骚扰的沉默文化、
以及制造出预备性特权（\secref{s:classroom-mixed}）的日益严重的社会不平等，
都不是任何一个人的错，
但修好它们是每个人的责任。
作为老师，
你比大多数人有更多力量；
\hreffoot{https://frameshiftconsulting.com/ally-skills-workshop/}{这个工作坊}
就如何做一个好的盟友给出了极好的实用建议，
而它的建议，大概比本书教给你的任何教学知识都更重要。

\seclbl{练习}{s:motivation-exercises}

\exercise{真实任务}{两人一组}{15}

\begin{enumerate}

\item
  两人一组，
  列出你这周做过的、用到了你所教技能的六件事。

\item
  把这些条目放到「掌握所需时间」与「有用程度」的 2x2 网格上。
  你们在哪里一致、在哪里分歧？

\end{enumerate}

\exercise{核心需求}{全班}{10}

Paloma Medina 识别出了职场中人的\hreffoot{https://www.palomamedina.com/biceps}{六种核心需求}：
归属、
改进（也就是取得进展）、
选择、
平等、
可预测、
和重要性。
读过她对这些需求的描述之后，
按对你个人的重要程度从高到低排序，
然后和同伴的排序比较。
你认为你的排序和你的学习者的排序相比如何？

\exercise{落实一项包容性策略}{个人}{5}

从~\cite{Lee2017}里挑一个你想着手做的活动或做法改变。
在三个月后的日历上设一个提醒，
到那时问自己有没有为此做点什么。

\exercise{事后复盘}{思考—结对—分享}{20}

\begin{enumerate}
\item
  回想你过去上过的一门课，
  找出老师做过的一件让你失去动力的事。
  记下事后本可以怎么做来纠正那种情况。
\item
  和邻座结对，比较你们的故事，
  然后把你的评论加到全班共享的一份笔记里。
\item
  作为一个群体，回顾共享笔记里的评论。
  挑出几件本可以做得不同的事，加以高亮和讨论。
\item
  你觉得这么做能帮你将来处理类似情况吗？
\end{enumerate}

\exercise{走一遍路线}{全班}{15}

找到离你们楼最近的公共交通下车点，
从那里走到你的办公室，再走到最近的洗手间，
记下你认为对行动不便者来说困难的地方。
现在借一把轮椅，重走一遍这段路。
你的清单有多完整？
又有没注意到，这个练习的第一句话本身就假设了你能走路？

\exercise{谁说了算？}{全班}{15}

在~\cite{Litt2004}中，
Kenneth Wesson 写道：
「如果贫困城区的孩子在社会标准化考试上一直压过富裕郊区的孩子，
还有谁会天真到相信我们仍然会坚持用这些考试作为成功的指标？」
读一读 Cameron Cottrill 的\hreffoot{https://mobile.nytimes.com/2016/04/10/upshot/why-talented-black-and-hispanic-students-can-go-undiscovered.html}{这篇文章}，
然后从你自己的经历里，描述一个「客观」评估反而巩固了现状的例子。

\exercise{常见的刻板印象}{两人一组}{10}

有些人还会说：「这太简单了，连你奶奶都会用。」
两人一组，
再列出两三个会强化计算领域刻板印象的说法。

\exercise{不当混蛋}{个人}{15}

Gary Bernhardt 的\hreffoot{https://www.destroyallsoftware.com/blog/2018/a-case-study-in-not-being-a-jerk-in-open-source}{这篇短文}
把一条不必要的敌意留言改写得不那么粗鲁。
以它为范本，
在 \hreffoot{https://stackoverflow.com/}{Stack Overflow} 或其他公开讨论区里找一条令人不快的留言，
把它改写得更具包容性。

\exercise{留面子}{个人}{10}

你希望招来的学习者里，有没有人会因为承认自己还不懂你想教的某些东西而难堪？
如果有，
你怎样帮他们保住面子？

\exercise{童年玩具}{全班}{15}

\cite{Cutt2017}调查了成年电脑用户童年的活动，
发现信心与电脑使用之间最强的相关，
建立在独自阅读、以及玩乐高这类没有活动部件的搭建类玩具上。
在全班做一次调查，看看人们还参与过哪些别的活动，
然后到网上去搜这些活动。
它们的介绍和广告有多强的性别色彩？
你认为这带来了什么影响？

\exercise{课程无障碍}{两人一组}{30}

两人一组，
选一门材料可以在网上找到的课程，
各自按照
\hreffoot{https://accessibility.blog.gov.uk/2016/09/02/dos-and-donts-on-designing-for-accessibility/}{这些海报}里的该做与不该做来给它打分。
你和同伴在哪里一致？
在哪里分歧？
这门课对六类用户的效果分别如何？

\exercise{追踪这个循环}{小组}{15}

\cite{Coco2018}追踪了一个令人沮丧地常见的模式：
良好的意愿，被一个组织的领导层不愿真正改变所破坏。
以 4—6 人为一组，
写出你想象中这个循环的每一阶段、
各相关方会发给对方的简短文字或邮件。

\exercise{最糟能糟到哪儿去？}{小组}{5}

这些年里，
我遇到过投影仪着火、
一名学生临产、
以及课堂上打起来。
我从台上摔下来过两次，
在自己的讲座上睡着过，
还讲砸过很多笑话。
以小组为单位，
列出你们讲课时遇到过的最糟的事，
然后分享给全班。
把这份清单留着，日后提醒自己：不管课有多糟，
至少\emph{那些}事没发生。

\section*{回顾}

\figpdfhere{figures/conceptmap-motivation.pdf}{概念：动机}{f:motivation-concepts}{}
